\section{Background}
\label{sec:background}

\subsection{Modular arithmetic as a group-theoretic task}

Let $p$ be an odd prime and let $\Z/p\Z = \{0, 1, \ldots, p-1\}$.
Under addition modulo $p$, $\Z/p\Z$ is a cyclic group of order $p$.
Under multiplication modulo $p$, the non-zero residues
$(\Z/p\Z)^* = \{1, 2, \ldots, p-1\}$ form a cyclic group of order $p - 1$
generated by some primitive root $g$. The additive and multiplicative
groups are both cyclic, but they have \emph{different} order, so their
Fourier bases live on different index sets.

\subsection{The discrete Fourier transform on $\Z/p\Z$}

For a function $f : \Z/p\Z \to \R$, the discrete Fourier transform (DFT) is
\begin{equation}
\hat f(k)
= \frac{1}{p} \sum_{n = 0}^{p - 1} f(n) \, e^{-2 \pi i k n / p}
\qquad k = 0, 1, \ldots, p - 1.
\label{eq:dft}
\end{equation}
The inverse transform reconstructs $f$ as a sum of complex exponentials
$e^{2\pi i k n / p}$, equivalently as a sum of cosines and sines indexed by
the frequency $k$. The DFT extends column-wise to matrices: a
$p \times d$ matrix $W$ has DFT $\widehat W \in \C^{p \times d}$ whose
$k$-th row records the contribution of frequency $k$ to each of the $d$
columns. The total power at frequency $k$ is
$\| \widehat W(k, :) \|_2^2$.

\subsection{The Fourier-multiplication algorithm for modular addition}

\citet{nanda2023progress} showed that the embeddings of a one-layer
transformer trained to grok $(a + b) \bmod p$ concentrate on a small set of
Fourier frequencies $\{k_1, \ldots, k_K\}$ with $K \approx 5$. The
algorithm the network has discovered is a multiplicative implementation of
the trigonometric product-to-sum identity
\begin{equation}
\cos(\alpha) \cos(\beta) - \sin(\alpha) \sin(\beta) = \cos(\alpha + \beta).
\label{eq:trig}
\end{equation}
Setting $\alpha = 2\pi k a / p$ and $\beta = 2\pi k b / p$, equation
\eqref{eq:trig} gives a way to compute features of $a + b \bmod p$ from
features of $a$ and $b$ multiplicatively --- exactly the kind of
nonlinearity that an attention layer followed by an MLP can implement.
The computation proceeds in three stages: the embedding matrix maps each
digit to its sine and cosine values at each of the $K$ frequencies; the
attention layer mixes the $a$ and $b$ representations across the residual
stream; and the MLP applies the product-to-sum identity. Reading out at
the final position recovers the result digit.

For \emph{multiplication} modulo a prime, the analogous algorithm uses a
character of the multiplicative group $(\Z/p\Z)^*$. Picking a primitive
root $g$ and writing each non-zero residue $a = g^{i_a}$, the
multiplicative DFT
\begin{equation}
\widetilde f(m) = \frac{1}{p - 1} \sum_{i = 0}^{p - 2}
                    f(g^i) \, e^{-2 \pi i m i / (p-1)}
\end{equation}
indexes by $m \in \{0, 1, \ldots, p-2\}$ rather than $\{0, \ldots, p-1\}$,
and the analogue of \eqref{eq:trig} is
$\cos(2\pi m i_a / (p-1)) \cos(2\pi m i_b / (p-1))
- \sin(2\pi m i_a / (p-1)) \sin(2\pi m i_b / (p-1))
= \cos(2\pi m (i_a + i_b) / (p-1))$,
which equals the same character evaluated at $i_a + i_b \bmod (p - 1)$.
That is the discrete-log of $a \cdot b \bmod p$. Hence a network that
solves $(a \cdot b) \bmod p$ via Fourier features must work on the
$(p - 1)$-cycle generated by $g$, not on the $p$-cycle of $\Z/p\Z$.

\subsection{The role of weight decay}

\citet{loshchilov2019adamw} introduced AdamW, which applies weight decay
directly to the parameter update rather than via the gradient. Standard
transformer training uses $\mathrm{weight\_decay} \in [0.01, 0.1]$. For
grokking the load-bearing setting is much larger:
$\mathrm{weight\_decay} = 1.0$. Without aggressive weight decay the
network sits in the memorisation regime indefinitely; the implicit
regularisation that drives the transition from memorisation circuits to
the Fourier algorithm requires a non-trivial preference for
small-Frobenius-norm solutions
\citep{liu2022understanding,liu2022omnigrok}.
